\section{Method}
\label{sec:method}

\subsection{Reference architecture}

The reference network is the source network of this study: a depth-$L$
fully-connected classifier with $L$ linear layers, the first $L-1$ of which
are followed by the same activation--normalization block, and the last of
which produces class logits without either. Writing $d_0$ for the input
dimension (784 for MNIST), $H$ for the hidden width, and $C$ for the number of
classes ($10$ for MNIST), the network is
\begin{align}
h^{(0)} &= x, \\
h^{(\ell)} &= \pi\bigl(f(W_{\ell} h^{(\ell-1)} + b_{\ell})\bigr),
  \qquad \ell = 1,\dots,L-1, \label{eq:block}\\
\hat{y} &= W_{L} h^{(L-1)} + b_{L},
\end{align}
with $x \in \mathbb{R}^{d_0}$ flattened from the input, $W_1 \in
\mathbb{R}^{H \times d_0}$, $W_\ell \in \mathbb{R}^{H \times H}$ for
$2 \le \ell \le L-1$, and $W_L \in \mathbb{R}^{C \times H}$. All weights are
initialized with Xavier-uniform and all biases to zero. The block
$h \mapsto \pi(f(Wh+b))$ is the object of study; both $f$ and $\pi$ are
independently selectable in our implementation, and the reference
configuration is $f(z) = z + \lambda z^{3}$ with $\lambda = 10^{-2}$,
$\pi(y) = y / \lVert y \rVert_2$.

\subsection{The projective polynomial block}

The reference activation is
\begin{equation}
f(z) = \operatorname{clamp}(z, [-c, c]) + \lambda \operatorname{clamp}(z, [-c, c])^{3},
\qquad c = 10,
\label{eq:poly}
\end{equation}
applied elementwise, followed by the per-sample projection
\begin{equation}
\pi(y) = \frac{y}{\max\{\lVert y \rVert_2, \epsilon\}},
\qquad \epsilon = 10^{-8},
\label{eq:proj}
\end{equation}
where the norm runs over the feature axis, so that $\pi(y)$ is a point of the
unit sphere and the layer map is invariant to the radial scale of $y$. We
emphasize three properties that motivate this choice and that our ablations
probe separately.
\begin{itemize}
\item \emph{Bounded growth.} The clamp bounds $|f(z)| \le c + \lambda c^3$,
  so the forward signal cannot blow up; the projection then bounds the
  \emph{norm} of each hidden activation at exactly $1$.
\item \emph{No statistics, no state.} $\pi$ has no parameters and no running
  buffers. Unlike BatchNorm it does not couple samples, and unlike both
  BatchNorm and LayerNorm it has no learned affine (`scale and shift'')
  parameters, so it changes the geometry of the representation without adding
  a single degree of freedom.
\item \emph{Radial invariance.} Because $\pi$ identifies $y$ with $y/\lVert y
  \rVert$, the block discards the radial direction at every layer, which is
  the sense in which the construction is projective.
\end{itemize}

We separate the two components in the implementation: \texttt{--activation
poly} gives Eq.~\eqref{eq:poly} \emph{without} the projection, \texttt{--norm
l2} gives Eq.~\eqref{eq:proj} \emph{without} the polynomial, and \texttt{--activation
poly\_proj} gives the reference block of Eq.~\eqref{eq:projective-poly}. This
factorization is what allows the same harness to attribute performance to the
polynomial, to the projection, or to their combination.

\paragraph{Placement.} Normalization can be applied before or after the
nonlinearity. Write the block as $\pi \circ f$ (``post'') or $f \circ \pi$
(``pre''). The reference construction is post-normalization; we report both,
since pre-normalization is standard in transformers while post-normalization
was the original BatchNorm convention.

\subsection{Activation families}

Let $z$ denote the pre-activation. All pointwise activations are applied
elementwise.

\paragraph{Pointwise.}
Identity (a linear control), the classical saturating $\tanh(z)$, and the
three non-saturating activations that current large language models use in
practice: $\operatorname{ReLU}(z) = \max(0,z)$, GELU $z\Phi(z)$ with $\Phi$ the
standard normal CDF (``exact'', not the tanh approximation), and SiLU
$z\sigma(z)$. We deliberately restrict the pointwise family to these, rather
than sweeping the many historical variants (LeakyReLU, ELU, SELU, Softplus,
Mish, sinusoidal), because none of those appears in a current frontier
language model; the study is meant to state where a polynomial activation
stands relative to what is actually deployed, not to re-audit the
activation-function zoo. The identity map is retained only as a control that
makes the contribution of nonlinearity itself measurable.

\paragraph{Polynomial.}
The fixed truncated cubic of Eq.~\eqref{eq:poly}, with the order and
coefficient treated as hyperparameters, and three \emph{learnable} polynomial
activations:
\begin{align}
f_{\text{Leg}}(z) &= g \sum_{k=0}^{K} a_k L_k\!\left(\frac{\tilde{z}}{B}\right),
\qquad
f_{\text{Cheb}}(z) = g \sum_{k=0}^{K} a_k T_k\!\left(\frac{\tilde{z}}{B}\right),
\qquad
\tilde{z} = \operatorname{clamp}\!\left(\frac{z}{s}, [-B, B]\right),
\end{align}
where $L_k$ and $T_k$ are the Legendre and Chebyshev polynomials of degree
$k$, $s$ and $B$ fix the normalization window, and $(a_k)$ together with the
gain $g$ are learned. We initialize $a_1 = 1$, $a_{k \neq 1} = 0$, $g = 1$,
so every learnable activation \emph{starts as the identity map} and must
justify any departure from it.
For the spline variant we use the KAN construction
\begin{equation}
f_{\text{spline}}(z) = w_{\text{base}}\,\operatorname{SiLU}(z) +
  w_{\text{spline}} \sum_{j} c_j N_j(\operatorname{clamp}(z,[-B,B])),
\end{equation}
with $N_j$ cubic B-spline basis functions on a uniform grid and control
points $c_j$ initialized to zero, so the activation also starts at
$w_{\text{base}}\operatorname{SiLU}(z)$ --- the KAN initialization.

\paragraph{Gated units.}
A gated unit replaces the pointwise map by
\begin{equation}
\operatorname{GLU}_{\sigma}(W x) = (W_v x) \odot \sigma(W_g x),
\label{eq:glu}
\end{equation}
with two matrices $W_v, W_g$ and a gate nonlinearity $\sigma$. We use the two
gates that current language models deploy --- $\sigma = \mathrm{SiLU}$
(SwiGLU, as in Llama and PaLM) and $\sigma = \mathrm{GELU}$ (GeGLU) --- and,
as the paper's polynomial counterpart, $\sigma = z + \lambda z^{3}$
(PolyGLU). Gated units with historical gate nonlinearities (ReGLU, SELU-GLU,
Mish-GLU, Tanh-GLU) and the fully degenerate \emph{bilinear} unit ($\sigma =
\operatorname{Id}$) are omitted for the same reason as the historical
pointwise activations: they are not what is deployed, and including them
would broaden the failure-mode survey without changing the comparison the
paper is about.

Because Eq.~\eqref{eq:glu} carries about twice the parameters of a pointwise
block at the same width, a pointwise-versus-gated comparison at equal width
confounds the activation with capacity. We therefore report gated results at
a width chosen so that the \emph{total} parameter count matches the pointwise
network, and we additionally report pointwise networks at that reduced width
as a control. Concretely, for a depth-$L$ network with input dimension
$d_0$ and $C$ classes, a pointwise network at width $W_e$ has
$d_0 W_e + (L-2)W_e^2 + C W_e + C$ weights, while a gated network at width
$W_g$ has $2d_0 W_g + 2(L-2)W_g^2 + C W_g + C$; we choose $W_g = 82$ to match
the pointwise width $W_e = 128$ as closely as an integer width allows at
$L = 10$, $d_0 = 784$, $C = 10$. At these settings the two blocks have
\paramElem{} and \paramGate{} parameters respectively, a residual gap of
\paramGapPct{}\%; \emph{no integer width removes this gap}, because the
gated block carries a second $d_0 W_g$ term and a second $(L-2)W_g$ multiplier,
which is not an affine function of the width. We report this gap rather than
hiding it, and we also report pointwise networks at the reduced width
$W_g = 82$ as a control, so that the activation effect can be read off the
matched-capacity rows as well as the matched-width rows.

\subsection{Normalization families}

We compare five normalization choices, all applied to the hidden activation
with the feature axis as the normalization axis.
\begin{itemize}
\item \emph{None}, the unnormalized control.
\item \emph{Projection} ($\pi$): Eq.~\eqref{eq:proj}, no parameters.
\item \emph{BatchNorm}: per-feature standardization by mini-batch mean and
  variance, with learned scale and shift and running statistics for
  evaluation.
\item \emph{LayerNorm}: per-sample standardization over the feature axis, with
  learned scale and shift.
\item \emph{RMSNorm}: per-sample root-mean-square scaling, with a learned
  scale but no mean subtraction and no shift.
\end{itemize}
This choice makes the comparison informative in a specific way: BatchNorm and
LayerNorm both add a learned affine map after normalizing, while projection
and RMSNorm (essentially) do not, so we can separate the effect of
\emph{normalizing} from the effect of \emph{re-parameterizing with a learned
affine map}.

\subsection{Diagnostics}

For each run we record, in addition to loss and accuracy curves, the
per-layer gradient norms and the per-layer activation norms at
initialization and after training, on a fixed mini-batch. These let us
distinguish runs that differ in accuracy because of optimization (gradient
profile) from runs that differ because of representation scale (activation
profile), and in particular they make visible the collapse of the
unnormalized gated unit reported in Sec.~\ref{sec:results}.
